\documentclass[11pt]{article}

\usepackage[T1]{fontenc}
\usepackage{amsmath,amssymb,amsthm}
\usepackage{booktabs}
\usepackage{array}
\usepackage{enumitem}
\usepackage{geometry}
\usepackage[hidelinks]{hyperref}
\usepackage{microtype}

\geometry{margin=1in}
\setlist{nosep}

\newtheorem{definition}{Definition}[section]
\newtheorem{proposition}[definition]{Proposition}
\newtheorem{remark}[definition]{Remark}

\newcommand{\bits}[1]{\lvert #1\rvert}
\newcommand{\pair}[2]{\langle #1,#2\rangle}
\newcommand{\K}{K}
\newcommand{\Eone}{E_1}
\newcommand{\Etwo}{E_2}
\newcolumntype{L}[1]{>{\raggedright\arraybackslash}p{#1}}

\title{Algorithmic Containment:\\
Directionality, Uniformity, and a Self-Delimiting Repair}
\author{MeTTapedia Formalization Project}
\date{August 2026}

\begin{document}
\maketitle

\begin{abstract}
Franz's proposal to understand implication through containment between
generative descriptions gives a promising non-probabilistic world-model
semantics for inference.  We report a source-level and machine-checked audit
of the two principal results in the 2026 formulation.  The displayed proof of
Theorem~1 reverses a directional mutual-information premise and later uses a
feature-compression hypothesis absent from its theorem statement.  The
directional step can be repaired with an explicit symmetry-of-information
loss, while the compression premise should be stated.  The construction in
Theorem~2 asks a program to recover a feature from
$\pair{g^*}{\K(g^*)}$, but its encoded parameter transports the program and
residual without transporting the exact value $\K(g^*)$.  We give a uniform
self-delimiting repair carrying the triple $(\K(g^*),t,p)$ and prove its
extraction, reconstruction, injectivity, and finite slack ledger in Lean.

These findings invalidate the two displayed proof routes, not the research
programme.  The corrected construction establishes exact image inclusion and
additive-slack feature transfer.  Completing the fully universal prefix-
complexity theorem still requires an effective conditional prefix machine,
the strong conditional chain rule, and symmetry of information with uniform
constants.  We state that boundary explicitly and invite correction of any
source interpretation in this preliminary note.
\end{abstract}

\section{Scope, conventions, and attribution}

The theory audited here is due to Arthur Franz, especially
\emph{Grounded Reasoning: Implication as Algorithmic Containment}
\cite{franz2026}, building on the incremental-compression theory of Franz,
Antonenko, and Soletskyi \cite{franz2021} and the earlier programme
\cite{franz2016}.  The prefix-complexity chain rule and symmetry results are
standard algorithmic-information theory; our reference point is
Li and Vit\'anyi \cite{livitanyi2019}.

The following distinctions are essential.

\begin{itemize}
  \item $\K(a\mid b)$ is prefix Kolmogorov complexity relative to one fixed
        universal prefix machine.
  \item The 2026 paper defines directional information by
        $I(a:b)=\K(a)-\K(a\mid b)$.  It is not definitionally symmetric.
  \item An $O(1)$ or $O(\log n)$ statement is meaningful here only as a
        uniform family statement: its hidden constant cannot depend on the
        particular strings being bounded.
  \item Exact featurehood, requiring
        $\bits f+\bits p<\bits x$, is different from featurehood with an
        additive logarithmic slack.  Our formalization gives them different
        types.
\end{itemize}

This note distinguishes three judgments: a displayed proof may be invalid; a
theorem statement may require an omitted premise; and an existential theorem
may remain true by a repaired construction.  We do not infer the last judgment
from either of the first two.

\section{Theorem 1: the directional step}

Theorem~1 of \cite{franz2026} assumes, in the paper's notation,
\begin{align}
  f(p)&=x,\\
  \K(g\mid x)&=O(\log\bits x),\\
  I(p:f,g)&=O(1),
\end{align}
and concludes $\K(g\mid f)=O(\log\bits x)$.

With the paper's directional definition, its third premise expands as
\[
  \K(p)-\K(p\mid f,g)=O(1).
\]
Equation~(16) of the proof instead uses
\[
  \K(f,g\mid p)=\K(f,g)+O(1).
\]
That is the reverse orientation.  It does not follow by unfolding the stated
definition.  A finite conditional prefix-machine example in the formal audit
realizes the two directional values as exactly zero and one bit, so the swap is
not a definitional equality.

The appropriate repair is not to declare mutual information exactly
symmetric.  Rather, one introduces the explicit symmetry ledger
\begin{equation}
  \K(f,g)+\K(p\mid f,g)
  \leq
  \K(p)+\K(f,g\mid p)+s_{\mathrm{sym}},
  \label{eq:symledger}
\end{equation}
where $s_{\mathrm{sym}}$ is logarithmic under the usual size controls.
Combining the source-oriented genericity inequality
\[
  \K(p)\leq\K(p\mid f,g)+s_{\mathrm{gen}}
\]
with \eqref{eq:symledger} gives the orientation actually used later:
\[
  \K(f,g)\leq\K(f,g\mid p)+s_{\mathrm{gen}}+s_{\mathrm{sym}}.
\]
This cancellation is a theorem over natural numbers; no signed subtraction is
needed in the proof ledger.

\subsection{The omitted compression premise}

The displayed assumptions of Theorem~1 do not include
\begin{equation}
  \bits f+\bits p<\bits x.
  \label{eq:featurecompression}
\end{equation}
The proof nevertheless invokes \eqref{eq:featurecompression} after its
equation~(19) to infer $\bits f<\bits x$ and hence to bound the cost of
transporting $\K(f)$.  The premise belongs in the theorem statement.  It is
available in Corollary~1 because that corollary assumes that $f$ is a feature
of $x$, but it is absent from the more general theorem as displayed.

\subsection{Uniform corrected form}

The machine-checked repair proves the following finite statement.  Assume:

\begin{enumerate}
  \item explicit programs for $f(p)=x$ and for recovering $g$ from $x$;
  \item representability of every complexity term used in a shortest-program
        argument;
  \item fixed, condition-uniform simulations of application, projection,
        literal, indexing, and pairing machines;
  \item a strong conditional chain rule with one uniform additive constant;
  \item the direct genericity inequality, or the source inequality together
        with an explicit symmetry ledger; and
  \item the compression premise \eqref{eq:featurecompression}.
\end{enumerate}

Then
\[
  \K(g\mid f)
  \leq \K(g\mid x)+2\log_2(\bits x+1)+C,
\]
where $C$ is displayed as a sum of the fixed compiler, chain-rule,
genericity, and symmetry slacks.  The formal proof uses only natural-number
inequalities.  The strong chain rule and symmetry theorem are currently
interfaces rather than constructed instances; consequently this is a proved
reduction of Theorem~1 to standard effective AIT infrastructure, not yet a
closed construction of that infrastructure.

A subtle uniformity point matters here.  A machine that hard-wires particular
conditions can establish an inequality for those particular strings, but its
compiler length may grow with those strings.  Such a machine cannot witness a
family-uniform $O(1)$.  The corrected formal proof therefore uses fixed
projection machines whose code is independent of $f,p,g,$ and $x$.

\section{Theorem 2: the missing complexity index}

Theorem~2 chooses a shortest program $t$ witnessing
\[
  \bits t=\K(f\mid g^*,\K(g^*)).
\]
Thus $t$ expects both $g^*$ and the exact integer $\K(g^*)$ as auxiliary
input.  The construction then sets
\[
  q=\Etwo(t)p
\]
and says that a fixed interpreter reconstructs $f$ from $(g^*,t)$.  The
encoded parameter contains $t$ and $p$, but not $\K(g^*)$.  In general the
missing integer is not computable from $g^*$; supplying it by a program whose
length is uniform in $g^*$ is precisely what cannot be assumed.

This is a gap in the displayed construction.  It does not by itself disprove
the existential conclusion of Theorem~2.

\subsection{A self-delimiting repair}

Let $k=\K(g^*)$, and encode
\[
  q'=\Etwo\bigl(\Etwo(\operatorname{bin}(k))t\bigr)p.
\]
Equivalently, $q'$ is a nested self-delimiting encoding of the triple
$(k,t,p)$.  A single fixed wrapper now:

\begin{enumerate}
  \item decodes $k,t,p$ from $q'$;
  \item runs $t$ with condition $\pair{g^*}{\operatorname{bin}(k)}$ to obtain
        $f$; and
  \item runs $f$ on $p$.
\end{enumerate}

No complexity oracle is used.  The wrapper is uniform in all four strings.
The corresponding descriptive program first uses the old descriptive program
to recover $p$ from the source and then emits the same triple $(k,t,p)$.
Hence extraction and reconstruction are both executable, rather than merely
an image-membership assertion.

The exact residual length is
\begin{align}
  \bits{q'}={}&
  \bits{\Etwo(\operatorname{bin}(k))t}
  +2\bits{\operatorname{bin}
      (\bits{\Etwo(\operatorname{bin}(k))t})}+1
  +\bits p.
  \label{eq:qprimelength}
\end{align}
The formalization proves this identity and the logarithmic upper bound.  It
also proves the necessary negative control: the map that transports only
$(t,p)$ is not injective as an encoding of triples $(k,t,p)$, since it erases
the first coordinate.

\subsection{What is proved, and what remains}

Let $G(g^*)$ be the tagged wrapper generator and let
$H(k,t)=\Etwo(\operatorname{bin}(k))t$.  Suppose the finite
description-transfer inequality
\begin{equation}
  \bits{G(g^*)}+\bits{H(k,t)}
  \leq \bits f+s_{\mathrm{transfer}}
  \label{eq:transfer}
\end{equation}
holds.  The executable Lean theorem proves that every strict $f$-feature
witness $f(p)=y$ becomes a $G(g^*)$-feature witness with slack
\[
  s_{\mathrm{transfer}}
  +2\bits{\operatorname{bin}(\bits{H(k,t)})}+1.
\]
The same fixed generator works for every such $p$.  Therefore the repaired
construction proves exact image inclusion and a universal additive-slack
feature implication.  A concrete finite positive instance is included.

Two boundaries should not be blurred.

\begin{itemize}
  \item The result is additive-slack featurehood, not the strict featurehood
        defined by $\bits g+\bits q<\bits y$.  The earlier theory already
        notes a logarithmic slack, and the formal types preserve that
        distinction.
  \item Deriving \eqref{eq:transfer} from the source's prefix-complexity
        premises still depends on the strong chain rule and the shortest-
        feature bounds.  The 2021 incompressibility theorem supplies the
        one-sided estimate $\bits{g^*}\leq\K(g^*)+O(1)$ needed in that ledger.
        The equality written as equation~(29) in the 2026 proof is stronger
        than the cited theorem, but the stronger direction is unnecessary.
\end{itemize}

\section{Machine-checked artifact map}

The audit is organized by concept rather than by author or paper:

\begin{center}
\begin{tabular}{L{0.38\linewidth}L{0.53\linewidth}}
\toprule
Module & Role \\
\midrule
\path{KolmogorovComplexity.SelfDelimitingCode} &
$\Eone/\Etwo$ round trips, prefix-free ranges, exact and logarithmic lengths,
and a negative comparison with unary headers. \\
\path{KolmogorovComplexity.DirectionalInformation} &
Finite machine separating the two directional information values. \\
\path{KolmogorovComplexity.ContainmentRepair} &
Fixed machine transformations, explicit Nat ledgers, and corrected direct-
and source-orientation forms of Theorem~1. \\
\path{WorldModel.ContainmentEncoding} &
Indexed triple residual, decoder round trip, injectivity, length bound, and
the index-erasure counterexample. \\
\path{WorldModel.FeatureContainmentConstruction} &
Executable wrapper, cross-algorithm implication, exact slack transfer, and a
nonvacuous finite positive instance. \\
\bottomrule
\end{tabular}
\end{center}

All listed crown theorems are free of admitted proofs and new axioms.  Their
reported dependencies are the ordinary logical foundations already used by
Lean and Mathlib.  The interfaces for the effective strong chain rule,
universal prefix simulation, and full symmetry theorem remain visibly
undischarged; no theorem in the artifact pretends otherwise.

\section{Assessment}

The source issues are substantive but localized.

\begin{center}
\begin{tabular}{L{0.29\linewidth}L{0.62\linewidth}}
\toprule
Claim & Audit status \\
\midrule
Theorem~1 proof as printed &
Invalid at equation~(16); it also uses an unstated compression premise. \\
Corrected Theorem~1 &
Reduced in Lean to fixed machine simulations, the strong conditional chain
rule, and an explicit symmetry bound. \\
Theorem~2 construction as printed &
Incomplete because the condition $\K(g^*)$ needed by $t$ is not transported. \\
Repaired image inclusion &
Proved operationally using the indexed residual $(\K(g^*),t,p)$. \\
Repaired compression transfer &
Proved with an exact finite slack from an explicit description-transfer
inequality; the remaining AIT derivation of that inequality is open. \\
\bottomrule
\end{tabular}
\end{center}

The strongest surviving idea is also the most useful for world-model
semantics: a generative premise and residual jointly produce a world, and a
consequence is model-level only when it is stable under admissible residual
variation rather than smuggled in by one residual.  The repaired formalization
makes that idea proof-relevant.  It therefore supports continued development
of algorithmic containment while separating recoverability, uniformity,
strict compression, logarithmic slack, and efficient discoverability.

\begin{thebibliography}{9}

\bibitem{franz2026}
A. Franz,
``Grounded Reasoning: Implication as Algorithmic Containment,''
in \emph{Artificial General Intelligence (AGI 2026)}, LNAI 16854,
pp. 238--252, 2026.
\href{https://doi.org/10.1007/978-3-032-33010-9_15}
{doi:10.1007/978-3-032-33010-9\_15}.

\bibitem{franz2021}
A. Franz, O. Antonenko, and R. Soletskyi,
``A Theory of Incremental Compression,''
\emph{Information Sciences}, vol. 547, pp. 28--48, 2021.
\href{https://doi.org/10.1016/j.ins.2020.08.035}
{doi:10.1016/j.ins.2020.08.035}.

\bibitem{franz2016}
A. Franz,
``Some Theorems on Incremental Compression,''
in \emph{Artificial General Intelligence (AGI 2016)}, pp. 74--83, 2016.

\bibitem{livitanyi2019}
M. Li and P. Vit\'anyi,
\emph{An Introduction to Kolmogorov Complexity and Its Applications},
4th ed., Springer, 2019.

\end{thebibliography}

\end{document}
